\part{Synthesis}

\chapter{What the project established, refuted, and left open}
\label{r:ch:synthesis}

\section{Claims and their strength}

{\small
\begin{longtable}{p{0.30\textwidth}p{0.17\textwidth}p{0.45\textwidth}}
\caption{Every substantive claim of the project and how strongly the evidence supports it.}\label{r:tab:claims}\\
\toprule
Claim & Strength & Evidence and boundary \\
\midrule
\endfirsthead
\toprule
Claim & Strength & Evidence and boundary \\
\midrule
\endhead
\bottomrule
\endfoot
Public cytology benchmarks hide a crop-to-parent hierarchy & established & Study~1: recovered and verified for four datasets; Herlev negative control at zero. \\
Random crop splits share parents and pixels & established & 92.6\% (SIPaKMeD), 99.97\% (CRIC) of test cells with a training parent; \RecOverlapSip\% / \RecOverlapCric\% / \RecOverlapCx\% with shared pixels. \\
The parent is recognisable & established & Study~1 probe 109$\times$ chance; Study~2 encoder \RecEncoderTimes$\times$ chance. \\
Crop-level splitting inflates accuracy & established (registered, powered) & H1 +\HoneEst~pp, all twelve architectures; H2 +\HtwoEst~pp; five inference methods agree. \\
The inflation is leakage, not diversity & established (registered, replicated) & H3 +\HthreeEst~pp, diversity \DivEst~pp; replicated on fresh seeds (+\HthreeRepEst~pp). \\
Exposure follows structure & supported & CRIC, with 29 cells per slide and high class--slide dependence, inflates ten times more than SIPaKMeD; predicted by the closed forms. \\
The leak needs no background & established (registered) & F3 +\FthreeEst~pp with every background masked; test-time removal leaves \RecBgKeep\%. \\
The leak is tied to the field's colour rendering & supported (registered), qualified & F2 removes \RecStainRemove\% at test time; F4 shows colour jitter in training does not prevent it. \\
A background context route carries the leak & not supported & H4 \HfourEst~pp (\HfourLo\ to \HfourHi). \\
Models apply a slide-level label prior (CRIC) & refuted & F1 reversed (\FoneEst~pp); equal gains for discordant and concordant cells. \\
The leak is memorised shared pixels & not supported (exploratory) & Equal advantage on non-overlapping crops. \\
The inflation is carried by linearly decodable field identity & refuted within its scope & F5, F7, F8: recoverability cut \RecoveryCutMin--\RecoveryCutMax\%, gap flat (pooled slope \FeightEst, CI \FeightLo\ to \FeightHi). \\
Saliency maps reveal the leak & not supported (descriptive) & Grad-CAM identical for leak and clean models. \\
Stronger grouping (P3) is more conservative than P2 & not supported & P2 $-$ P3 \PtwoPthreeEst~pp; the ordering passes for \RecOrderingPassed\ of 12 architectures. \\
\end{longtable}
}

\section{What we got out of it}

\begin{workedbox}{Scientific results}
A certified, powered estimate of how much crop-level splitting inflates two public cytology benchmarks; a causal
separation of leakage from diversity; a mechanism narrowed by elimination; and a methodological warning that the
standard audit tools (probes, saliency, augmentation) cannot detect or remove this kind of leakage.
\end{workedbox}

\begin{workedbox}{Reusable artifacts}
Verified crop-to-parent maps for SIPaKMeD, Cx22, CRIC, and Herlev; registered P1/P1$'$/P2/P3 splits with hashes for 20
seeds; the 2C factorial designs and transplant plans; a risk calculator (\texttt{cytofair risk-report}) that predicts a
benchmark's exposure from its parent table; a deterministic GPU-resident training stack with complete per-run artifacts;
a fail-closed scheduler and collectors; exported records of every run (\artifact{results/study2/run_records}); 121~GB of
archived predictions, features, and weights for future analyses without retraining.
\end{workedbox}

\begin{workedbox}{Writing}
A single-column manuscript (\artifact{paper/v2}, 17 pages, every number generated), the earlier ISBI draft, the Study~1
record, and this lifecycle record.
\end{workedbox}

\chapter{What worked and what did not}
\label{r:ch:worked}

{\small
\begin{longtable}{p{0.18\textwidth}p{0.38\textwidth}p{0.36\textwidth}}
\caption{What worked and what did not, by area.}\label{r:tab:worked}\\
\toprule
Area & Worked & Did not work, or cost more than expected \\
\midrule
\endfirsthead
\toprule
Area & Worked & Did not work, or cost more than expected \\
\midrule
\endhead
\bottomrule
\endfoot
Question design & Separating structure, availability, consequence, and mechanism kept every claim at its true strength. & Study~1 asked the consequence question six times and could not answer it. \\
Statistics & One crossed mixed model with a split random effect, sized by simulation; registered sensitivity analyses that all agreed; the assumption trigger worked as designed. & Per-architecture testing in Study~1; Satterthwaite $df$ for H3 is only 3 with four architectures, so its interval is wide. \\
Causal designs & The 2C factorial cleanly separated leakage from diversity; evaluation-only follow-ups on saved models were cheap and decisive. & The 2B transplant contrast was noisy (sibling backgrounds help even clean models); the 2A MLP adversary was evaded by ConvNeXt-T. \\
Mechanism & Elimination with several independent manipulations gave a coherent picture. & No single test identifies the carrier positively; F2 and F4 pull in different directions. \\
Engineering & GPU-resident augmentation (exact torchvision port); packed hashed caches; bitwise determinism; complete per-run artifacts; zero failed tasks in 2,921. & The four-core CPU quota had to be discovered the hard way; the slow network home forced a strict storage discipline. \\
Governance & Register, smoke-test on non-contrast splits, dry-run, collect fail-closed, amend with dates; every deviation recorded. & A dry run that bypassed one pipeline stage let a labelling bug through to the first real analysis (Amendment~9). Registrations are local commits, not yet externally timestamped. \\
Operations & Detached controllers, status files, sequential offloading with verification, clean vacating of node-local storage. & Several shell and remote-wait mistakes (self-matching process searches, hung SSH sessions) cost time but no data. \\
\end{longtable}
}

\section{Demotivations, and why they mattered}

Four results were disappointing at the time. Study~1's six tests failed to survive correction, which looked like a null;
Study~2 showed it was underpowered. H4 and H5 failed to isolate a mechanism; Amendment~8 showed why (the leak is neither
in the background nor in a linearly readable code). F1 went the opposite way from the prediction; that reversal is itself
evidence against a slide-level label prior. And the ConvNeXt-T adversary was evaded; that failure is a concrete example
of why decodability and use must be measured separately. Each setback narrowed the explanation instead of ending it, and
none was hidden.

\section{Lessons}

\begin{enumerate}[leftmargin=1.8em]
\item Ask a pooled question once, with the dependence structure in the model (the shared test set per seed is a random
effect, not twelve independent tests).
\item Keep every per-sample artifact. Half of Amendment~8 cost no training because Study~2 kept logits, features, and
weights.
\item Measure manipulations independently of the mechanism they target: the gradient-reversal adversary's own accuracy
fell while an independent probe's rose.
\item Dry runs must pass through every stage of the real pipeline.
\item On shared hosts, discover the real resource limits (CPU quota, writable disks, network storage behaviour) before
designing the execution.
\end{enumerate}

\chapter{Decision log}
\label{r:ch:decisions}

Every design or execution decision that shaped a result, with the reason it was taken and the alternative it displaced.

{\small
\begin{longtable}{p{0.22\textwidth}p{0.40\textwidth}p{0.30\textwidth}}
\caption{Decisions and their reasons.}\label{r:tab:decisions}\\
\toprule
Decision & Reason & Alternative not taken \\
\midrule
\endfirsthead
\toprule
Decision & Reason & Alternative not taken \\
\midrule
\endhead
\bottomrule
\endfoot
\multicolumn{3}{l}{\textbf{Study 2 statistics}} \\
Add a split random effect to H1 & Twelve architectures share each seed's test set; without it the model treats twelve evaluations of one test set as independent. & The roadmap's formula without it. \\
Twenty seeds & Smallest count meeting the registered power rule (0.90 at pilot variances, 0.75 at 1.5$\times$). & 15 seeds (failed the inflated criterion); 25--30 (more compute for little power). \\
Twelve architectures & Breadth across CNNs, hierarchical transformers, plain ViTs, and a self-supervised model, so the effect is not tied to one family. & The six Study~1 models only. \\
Bootstrap primary when residuals are heavy-tailed & Registered before data; protects the primary test from non-normal seed-level noise. & Always the Satterthwaite $t$-test. \\
H2--H5 in one Holm family; H1 alone & H1 is the question; the mechanism tests are a family of related but separate claims. & One family of five (would have penalised the primary). \\
\multicolumn{3}{l}{\textbf{Study 2 designs}} \\
220 crops per class in 2C & Largest budget feasible for every registered seed and 100/100 unrelated seeds; 240 failed for seed 307. & 240 or 260 crops. \\
Common test set across all four 2C cells & Every cell is scored on the same crops, so differences are training effects only. & Separate test sets per cell. \\
Diversity estimated within clean cells only & The low-diversity leak cell gives each test crop more siblings, so the interaction confounds dose with diversity. & A full 2$\times$2 interpretation. \\
Sibling background as the 2B control & Isolates the cell's own acquisition context from class-typical context. & Only unseen-field backgrounds. \\
Class-conditional parent adversary & Suppressing parent identity must never require removing class information. & An unconditional 424-way adversary. \\
CRIC as $100\times100$ windows & CRIC releases nucleus coordinates on slide images, not crops; fixed windows are the simplest faithful crops. & Variable-size crops from annotations. \\
\multicolumn{3}{l}{\textbf{Execution}} \\
GPU-resident augmentation & Four-core CPU quota starved data loaders; an exact port keeps the registered recipe. & PIL augmentation in DataLoader workers. \\
One training job per GPU & GPUs were saturated by one job; two jobs shared throughput without speed-up. & Two jobs per GPU. \\
Math attention kernel only & Fused attention kernels have non-deterministic backward passes. & Faster fused kernels. \\
Float16 checkpoints & Halves storage of about 117~GB of weights; re-evaluation reproduces 418/480 models exactly, the rest within three crops. & Float32 checkpoints. \\
Amendment 7 before any 2A run & Deterministic divergence would otherwise have failed the whole arm. & Handling divergence after seeing it. \\
Registered analysis run on the GPU host & Run records point at node-local artifacts; moving gigabytes to a workstation was slower than computing in place. & Pulling all artifacts to the workstation. \\
\multicolumn{3}{l}{\textbf{Amendment 8}} \\
A separately registered follow-up & Study~2's mechanism tests were inconclusive, and saved artifacts made decisive tests cheap; registering first keeps them confirmatory for new questions. & Unregistered exploration, or re-testing H4/H5 with more seeds. \\
Reuse saved models where possible & F2 and F5 needed no training; three arms cost 7~GPU-hours in total. & Retraining for every arm. \\
Fresh seeds for retraining arms & New data for new questions; no reuse of Study~2 mechanism seeds. & Reusing the mechanism seeds. \\
Probe-matched linear adversary & ConvNeXt-T had evaded the MLP adversary; a standardised linear head cannot hide identity from a linear probe while leaving it readable. & Re-running the MLP adversary. \\
LEACE ranks up to 256 & The smallest feature dimension is 384 (DINOv2 ViT-S); 256 removes most within-class parent structure in every backbone. & Full-rank erasure (exceeds some feature dimensions). \\
HED-strong colour augmentation & A published, strong, stain-shaped perturbation. & Hue/saturation jitter only. \\
One device probe for every arm & Recoverability must be measured identically across manipulations; agreement with sklearn verified ($r=\ProbeAgreement$). & Mixing probe implementations. \\
Pixel-overlap split labelled exploratory & Added after the registered results to test a specific alternative explanation. & Presenting it as confirmatory. \\
\multicolumn{3}{l}{\textbf{Close-out}} \\
Offload checkpoints during the follow-up matrix & 52~GB of new checkpoints against a 71~GB margin; verified copies before deletion. & Pausing the matrix on the disk guard. \\
Recommend ISBI framing around audit tools & The registered evidence is strong; the novel take-home for an imaging audience is that probes and saliency cannot detect this leakage. & Framing around ``leakage exists''. \\
\end{longtable}
}

\chapter{Validity, recommendations, and next steps}
\label{r:ch:validity}

\section{Threats to validity}

\paragraph{Parents are not patients.} SIPaKMeD fields, Cx22 source images, and CRIC slide images are the deepest released
units. Several can come from one patient, so parent-disjoint evaluation is necessary but may not be sufficient for
clinical independence; every inflation estimate is a lower bound on what patient-level grouping would remove.

\paragraph{One recipe.} Every model used one fixed fine-tuning recipe. Stronger augmentation, longer training, or other
optimisers could change the size of the inflation and the models' reliance on acquisition context.

\paragraph{Smaller training sets in the factorial.} Matching crop and field counts required 220 crops per class (1,100 in
total against 3,239 in Phase~1); the 2C, 2B, and carrier effects are measured on smaller training sets than the headline
contrast, and the decomposition that transfers them is approximate.

\paragraph{Imperfect interventions.} Foreground masks are threshold rules; inpainting leaves smoothed backgrounds; Reinhard
transfer matches only first and second moments; erasure and adversaries act on pooled features and linear structure.

\paragraph{Post-results follow-up.} Amendment~8 was motivated by Study~2's results. Its hypotheses and analyses were fixed
before any follow-up run, but they are follow-up questions; the pixel-overlap split was added afterwards and is
exploratory.

\paragraph{CRIC crops.} The CRIC cells are windows cut by this project around published nucleus coordinates, not an
official crop release, so absolute CRIC accuracies are not directly comparable with published numbers.

\paragraph{External verification.} The registration commits have not yet been pushed or deposited; until they are, their
dates rest on the repository and on the commit hashes recorded inside every run.

\paragraph{No clinical claim.} The project audits research benchmarks; it does not estimate screening sensitivity,
specificity, prevalence-calibrated performance, or effects on readers.

\section{Recommendations}
\label{r:ch:recommendations}

\begin{enumerate}[leftmargin=1.8em]
\item \textbf{Release the hierarchy.} Every observation should carry its acquisition parent (field, source image, or slide),
a patient or specimen identifier where permitted, and, if audited, a duplicate-component identifier.
\item \textbf{Split at the deepest released unit and publish the split}, with its hash.
\item \textbf{Report both units}: observation and parent counts per split, and the fraction of test observations with a
training sibling.
\item \textbf{Check the structure before training} with the closed forms (\texttt{cytofair risk-report}).
\item \textbf{Treat one question as one test}, with a mixed model that treats architectures and splits as random effects,
sized by simulation.
\item \textbf{Do not certify a split with a probe, a saliency map, or an augmentation.} Leakage inflated accuracy with field
identity erased from the features, with identical Grad-CAM maps, with backgrounds removed, and with strong colour
augmentation.
\item \textbf{Keep the evidence}: per-sample predictions, final weights, and the exact split for every run.
\end{enumerate}

\section{Publication options and open decisions}
\label{r:ch:publication}

\subsection{Venues}
\label{r:sec:publication}

\paragraph{ISBI 2027.} Four pages including figures and tables (references may use a paid fifth page); deadline Monday 26
October 2026, 23:59 US Eastern; single-blind; arXiv preprints allowed; a ``Compliance with Ethical Standards'' section and a
disclosure of AI-generated content in the acknowledgments are required; presentation in person in Lausanne. The evidence is
stronger than typical for the venue. The risks are perceived novelty (patient-level splitting is textbook advice) and the
page limit. The recommended framing is ``how much crop-level splitting inflates public cytology benchmarks, and why the
usual audit tools cannot see it'': H1--H3 as the result, one figure (per-architecture SIPaKMeD with CRIC), one table of all
thirteen registered tests, and a paragraph on F2, F3, and the probe and Grad-CAM findings. Everything else waits for the
journal version.

\paragraph{Journal version.} The full manuscript (\artifact{paper/v2}) with Study~2 and Amendment~8. Check the journal's
policy on extended conference papers if ISBI is pursued first; ISBI only forbids simultaneous submission of a substantially
similar paper to another conference or workshop.

\paragraph{Preprint.} Allowed by ISBI; it would also give the registrations a public timestamp once the repository is pushed.

\subsection{Decisions only the authors can make}

\begin{itemize}[leftmargin=1.6em]
\item Push the \texttt{study2} branch (and the Cx22 commits) to the public repository, and deposit the preregistrations with
a third party (for example OSF) for an external timestamp.
\item Choose a software licence; confirm redistribution terms for derived artifacts (SIPaKMeD, CCEDD).
\item Co-author approval of every posted or submitted version; fill the e-mail and repository placeholders.
\item Whether to submit to ISBI 2027, and who presents in person.
\item The AI-use disclosure wording: much of the Study~2 code, analysis, manuscript text, and this record was drafted with
an AI assistant under the authors' direction.
\item Whether to open Phase~3 (other imaging domains) and the direct test of the F2--F4 tension.
\end{itemize}

\subsection{Future work}

\begin{itemize}[leftmargin=1.6em]
\item Apply test-time stain transfer to the colour-augmented carrier models to resolve the F2--F4 tension directly.
\item Non-linear and spatial (pre-pooling) erasure of field identity.
\item Patient-level grouping where a dataset's provenance allows it.
\item Other recipes (stronger augmentation, longer schedules, stain normalisation as standard practice).
\item Other domains where crops share acquisition units (the pre-specified fish-imagery V2).
\end{itemize}
